% Ce chapitre repose sur :
%   DOC: s22, s33, s34, s35, s36, s37
%   DOC: s38, s39, s40, s41, s42, s43
%   DOC: s44, s45, s46, s47, s48, s49
%   DOC: s50
%   OBS: (aucun)
%   HYP: choix-composition, choix-conteneurs, choix-controles, choix-environnement-verrouille, choix-generation, choix-integration-continue, choix-langage-unique, decision-moteur-en-memoire, decision-orchestrateur, decision-rapprochement-probabiliste, decision-schema-du-tableau-de-bord, decision-serveur-de-base, decision-transformation-declarative
%   CHI: (aucun)
%   SER: (aucune)
%   SECTIONS: 2

% La colonne Outil est composée en drapeau à gauche : justifiée sur sa largeur étroite, elle
% coupait « Apache Air-flow » et « Docker Com-pose ». En drapeau, chaque nom d'outil tient sur sa
% ligne et n'a plus à être coupé.
%
% Chaque ligne des deux tableaux porte exactement une commande de citation, vers la documentation
% officielle de l'outil qu'elle nomme ; un contrôle le vérifie. Le motif d'un choix porte
% l'étiquette `decision-…` quand un enregistrement de décision du projet l'a consigné, et
% `choix-…` quand il est un choix posé sans enregistrement.

\chapter{Environnement et outils}\label{chap:outils}

La chaîne repose sur une dizaine d'outils libres. La question traitée est celle de leur choix : ce
que chacun est, ce que le projet en fait, et le motif qui l'a fait retenir. Un motif vient d'une
décision consignée par le projet, ou il est déclaré comme un choix posé.

\section{Les briques de la chaîne}

Chaque étape de la chaîne a son outil, du fichier déposé à l'écran (cf. tableau~\ref{tab:briques}).
Les transformations s'écrivent en SQL (\emph{Structured Query Language}), le langage de requête des
bases relationnelles \cite{s38}.

{\footnotesize
\begin{longtable}{@{}>{\raggedright\arraybackslash}p{2.6cm}p{3.4cm}p{4.0cm}p{4.6cm}@{}}
\caption[Les outils de la chaîne]{Les outils de la chaîne et le motif de leur choix}\label{tab:briques}\\
\toprule
Outil & Ce qu'il est & Ce que le projet en fait & Motif du choix \\
\midrule
\endfirsthead
\toprule
Outil & Ce qu'il est & Ce que le projet en fait & Motif du choix \\
\midrule
\endhead
Python\cite{s33} & langage de programmation & génère et charge les données, rapproche les fiches, sert le tableau de bord & un seul langage pour toute la chaîne\convention{choix-langage-unique} \\
\addlinespace
uv\cite{s34} & gestionnaire de dépendances Python & installe le même environnement partout, depuis un fichier de verrouillage & un environnement reproductible\convention{choix-environnement-verrouille} \\
\addlinespace
Faker, pandas et Arrow\cite{s35, s36, s37} & bibliothèques de valeurs fictives, de tables et de format colonnaire & produisent le jeu de données et l'écrivent en fichiers & des valeurs plausibles sans donnée réelle\convention{choix-generation} \\
\addlinespace
PostgreSQL\cite{s38} & serveur de base de données relationnelle & porte toutes les couches et l'instantané & plusieurs services partagent la base, et le chargement écrit pendant la lecture\convention{decision-serveur-de-base} \\
\addlinespace
dbt\cite{s39} & outil de transformation en SQL & type, modélise et teste chaque couche & des transformations écrites en SQL, testées couche par couche\convention{decision-transformation-declarative} \\
\addlinespace
Splink\cite{s22} & bibliothèque de rapprochement probabiliste & prédit les paires de fiches d'une même personne & un modèle estimé sur les données retrouve aussi les fiches qui diffèrent\convention{decision-rapprochement-probabiliste} \\
\addlinespace
DuckDB\cite{s40} & moteur de base de données en mémoire & exécute les calculs du rapprochement & ses tables de travail restent hors du serveur\convention{decision-moteur-en-memoire} \\
\addlinespace
Apache Airflow\cite{s41} & orchestrateur de tâches & lance chaque jour les tâches dans l'ordre de leurs dépendances & les tâches s'exécutent dans l'environnement du projet\convention{decision-orchestrateur} \\
\addlinespace
Streamlit\cite{s42} & bibliothèque d'application web & sert les pages du tableau de bord & un schéma dédié, remplacé en une transaction\convention{decision-schema-du-tableau-de-bord} \\
\addlinespace
Docker Compose\cite{s43} & outil de composition de conteneurs & démarre la base, l'orchestrateur et le tableau de bord & la même base sur le poste et en intégration continue\convention{choix-conteneurs} \\
\bottomrule
\end{longtable}
}

Le serveur relationnel montre le genre de motif retenu. Un moteur embarqué aurait produit les mêmes
tables plus vite, mais la chaîne a besoin que plusieurs services partagent la base par le réseau
(section~\ref{sec:choix-serveur}).

\section{Les outils de qualité et de composition}

D'autres outils contrôlent la chaîne et composent les documents (cf.
tableau~\ref{tab:outils-qualite}).

{\footnotesize
\begin{longtable}{@{}>{\raggedright\arraybackslash}p{2.6cm}p{3.4cm}p{4.0cm}p{4.6cm}@{}}
\caption[Les outils de vérification et de composition]{Les outils de vérification et de composition, regroupés par fonction}\label{tab:outils-qualite}\\
\toprule
Outils & Ce qu'ils sont & Ce que le projet en fait & Motif du choix \\
\midrule
\endfirsthead
\toprule
Outils & Ce qu'ils sont & Ce que le projet en fait & Motif du choix \\
\midrule
\endhead
pytest, Ruff et SQLFluff\cite{s44, s45, s46} & cadre de contrôles, analyseurs du code Python et du SQL & exécutent les contrôles et bloquent les écarts de forme & un défaut arrêté avant la publication\convention{choix-controles} \\
\addlinespace
GitHub Actions\cite{s47} & service d'intégration continue & rejoue les contrôles et compose le rapport & une vérification hors du poste de travail\convention{choix-integration-continue} \\
\addlinespace
LaTeX, biblatex et pgfplots\cite{s48, s49, s50} & composition, bibliographie et graphiques & composent le rapport et ses graphiques depuis les données du registre & un document recomposé à l'identique\convention{choix-composition} \\
\bottomrule
\end{longtable}
}

\section*{Synthèse}

La chaîne tient en un langage, une base relationnelle et quelques outils libres, installés à
l'identique partout. Chaque outil a été retenu pour une propriété que la chaîne exerce, et son
motif est consigné. Le chapitre~\ref{chap:activite} analyse l'activité que cette chaîne mesure.
